\section{Gemini 4 Argon、グローバル提供を開始}
Google DeepMindはGemini 4 Argonを開発者、企業、一般利用者向けにグローバル公開したと発表した。コーディング、企業の複雑な業務、サイバー防御などを想定し、Vals Index首位やDeepSWE v1.1で77.9\%といった評価値を示している。API導入価格は11月30日まで入力100万トークン当たり2ドル、出力10ドルと説明された。Antigravityでの利用も案内されている。これらの性能値は主に公式発表に基づき、独立した比較検証は本号では実施していない。
\dailyxsource{Google DeepMind (@GoogleDeepMind)}{https://x.com/GoogleDeepMind/status/2110039274922131792}

\section{OpenAI、社内録音機器の発見を報告}
OpenAI Newsroomは、オフィス内に隠された録音デバイスを発見し、AIを用いた調査で設置者を特定したと投稿した。投稿によれば、従業員による設置が判明して解雇に至り、FBIへの通報と逮捕もあったという。モデルの重みや顧客データの漏洩兆候はないと説明している。ただし、Grok収集では公式X投稿以外の一次資料や捜査記録は確認されていない。事件の詳細やAIによる分析方法は、現時点では同社の説明として受け取る必要がある。
\dailyxsource{OpenAI Newsroom (@OpenAINewsroom)}{https://x.com/OpenAINewsroom/status/2110095403198849132}

\section{Anthropic、独自ソース管理「Scion」を構築}
Anthropicは、エージェントを用いた開発でコードのマージが5月の1日約2,500件から9月には3万件超へ増え、従来のGit運用がボトルネックになったと報告した。対策として、サーバー側コミットやスパースチェックアウトを採用する常時接続型のソース管理基盤「Scion」を約5か月で構築・移行したという。インフラ規模を200分の1に縮小し、CPU使用量を25\%削減したとの効果も示された。数値は同社のエンジニアリング報告に基づく。
\dailyxsource{Anthropic (@AnthropicAI)}{https://x.com/AnthropicAI/status/2110047081414987957}

\section{Gemini 4 Argonのロボティクス実証}
Google DeepMindはGemini 4 Argonを使い、動画から物理的な状況を把握してシミュレーションを作成し、Apptronikの人型ロボットApolloにコーンホールを行わせるデモを公開した。7回連続成功したと説明し、物理理解、Real2Sim、反復的な改善を強調している。開発者向け記事ではERQA評価87\%などの数値も紹介された。ただし、デモの再現性や他モデルとの比較条件はGrokの収集範囲では確認されておらず、公開事例として紹介する。
\dailyxsource{Google DeepMind (@GoogleDeepMind)}{https://x.com/GoogleDeepMind/status/2110078384617464201}

\section{Liquid AIのMoEモデル、オープンウェイトで紹介}
Liquid AIの「LFM2.5-24B-A2B」が、オープンウェイトのオンデバイス向けモデルとして共有された。総パラメータ約240億、稼働時約20億のMoE構成、131Kコンテキスト、ツール呼び出し対応を掲げる。Apple M5 Maxで毎秒184トークンという速度も投稿で紹介されたが、Grokは公式ブログの全文を取得しておらず、実測条件やライセンスの詳細は未確認である。発表者・紹介者の数値を一般的な性能保証とは扱わない。
\dailyxsource{AgentPlays (@Agent\_PlaysAi)}{https://x.com/Agent_PlaysAi/status/2110128378300821588}

\section{Fish Audio「Drama 3」の音声評価が話題}
オープンウェイト音声モデル「Drama 3」について、Artificial Analysisの発音テストで93.2\%を達成し、従来の88.1\%から改善したとの投稿が共有された。年間収益50万ドル未満なら商用利用が無料という条件も紹介されている。ただし、今回のGrok収集ではFish Audioの一次発表や評価元の結果、ライセンス本文を確認できていない。順位・精度・利用条件はいずれも第三者投稿に基づく未検証情報として記録する。
\dailyxsource{The AI Post (@ThePost\_AI)}{https://x.com/ThePost_AI/status/2110128317718577384}

\section{MiniCPM-V 4.7の端末実行に注目}
OpenBMBのマルチモーダルモデル「MiniCPM-V 4.7」について、1Bの密なモデルと35B-A3BのMoE構成が再び紹介された。投稿では、iPhone 17 Pro Max上で毎秒13トークン超、メモリ使用量約3GBとする実行例や、llama.cpp、Ollama、vLLMへの対応が挙げられている。前日のモデル公開を受けた継続的な話題であり、本日の新発表とは区別する。端末性能や対応範囲は一次資料で確認できておらず、投稿者の主張として扱う。
\dailyxsource{Simon Chat (@SimonLe99917871)}{https://x.com/SimonLe99917871/status/2110126683361583553}

\section{Meta Muse Wonderをめぐる比較が継続}
Metaの「Muse Wonder」と上位版「Wonder Full」について、Gemini 4 Argonとの比較やエージェント用途での評価がX上で続いた。複数の投稿はActive Memoryなどの特徴に触れ、Metaがフロンティアモデル競争へ戻ったと評している。ただし、公式発表は今回の観測窓より前だった可能性が高く、各種ベンチマークや機能の細部も本号では確認していない。新製品公開の速報ではなく、既存発表への継続的な反応として記録する。
\dailyxsource{m0h (@exploraX\_)}{https://x.com/exploraX_/status/2110030971512766846}

\section{Anthropic、蒸留攻撃への防御改善を報告}
Anthropicは、Claudeに対する蒸留攻撃の有効性を大きく下げる防御策の進展を公式Xで報告した。観測窓開始直後の投稿であり、関連する研究記事へのリンクも示されている。二次投稿では攻撃の成功を90～99\%程度抑制したとの数値が言及されたが、Grok収集では防御方式や評価条件の詳細を確認していない。改善率を確定した一般的な効果として扱わず、モデルの不正な模倣への対策が議論された事例として紹介する。
\dailyxsource{Anthropic (@AnthropicAI)}{https://x.com/AnthropicAI/status/2109795898909950247}
